\section{Related work}
\label{sec:zflean-related}
We position our contributions with respect to set-theoretic developments in interactive provers such as Isabelle and Rocq, to Mathlib's ZFC model, and to categorical/set-theoretic bases.

\paragraph*{Isabelle/ZF and ZFC in Isabelle/HOL}
Isabelle includes a classical first-order logic (FOL) object logic with Zermelo–Fraenkel set theory as an object-level theory, commonly referred to as Isabelle/ZF~\cite{paulson-zf}.
In this approach, set-theoretic reasoning happens inside the ZF object logic, with membership and extensionality primitives and large portions of mathematics developed directly at the set level.
While powerful, ZF lives as a separate object logic from Isabelle's higher-order logic (HOL).
This separation historically supported extensive set-theoretic formalizations---from elementary theories such as ordinals up to genuinely deep developments such as the constructible universe and relative consistency results---but makes it less convenient to reuse HOL automation, type classes, and libraries in mixed developments.

A complementary line of work encodes ZFC \emph{within} Isabelle/HOL by introducing a HOL type of sets together with an elementhood relation and basic operations, with an emphasis on close integration with HOL's automation and type classes~\cite{paulson-zf-hol}.
In this setting, sets are first-class HOL values; proofs benefit from standard HOL tooling and can interoperate smoothly with typed libraries.
Conceptually, our Lean framework plays a similar bridging role: we keep ZFC-level objects and proofs, but we provide explicit bridges to Lean's native types so that tactics, typeclasses, and algebraic infrastructure can be reused where they pay off.
Our distinct contribution is a \emph{rewriting-friendly relational calculus} tailored to set-level developments, whereas the HOL encodings typically emphasize a minimal new notation surface and type-class integration.

\paragraph*{Rocq (ZFC, IZF)}
Rocq has multiple set-theoretic developments with different aims. One family (e.g. \texttt{coq-zfc}) encodes ZFC à la Aczel with primitive membership and extensional equality, and then \emph{proves} the ZFC axioms as theorems of the Rocq calculus~\cite{aczel-set-theory}.
Another strand (e.g. \texttt{coq-izf}, constructive or intuitionistic set theories) targets constructive metatheory or models of type theory, sometimes via pointed-graph semantics~\cite{miquel-pts-set-theory} or Tarski–Grothendieck universes~\cite{barras-set-theory}.
These developments showcase that strong set-level mathematics is possible inside a dependent type theory, but integration with Rocq's typed algebraic hierarchy and automation usually remains more ad-hoc.

\paragraph*{ETCS and categorical foundations}
Lawvere's Elementary Theory of the Category of Sets (ETCS) axiomatizes the (well-pointed) category of sets with finite limits, cartesian closure, a natural numbers object, and (typically) choice; equality is extensional at the \emph{morphism} level and ``up to isomorphism'' reasoning is built in from the start~\cite{lawvere-etcs}.
ETCS is weaker than ZFC and is rather a metatheoretic object than a framework for set-level development.
In principle though, ETCS could be pursued in Lean via category-theoretic libraries.

\medskip
Overall, our Lean framework shares with the above works the goal of enabling set-level mathematics inside a typed proof assistant, namely Lean, with smooth interoperability with typed libraries and automation.

